\section{Свойства беспроводного канала связи}
\label{sec:channel}

\subsection{Распространение сигналов и влияние земной поверхности}

Передача сигналов в беспроводных системах осуществляется с помощью радиоволн, которые излучаются антеннами в открытое пространство. Простейшей моделью распространения является модель свободного пространства. При изотропном излучении плотность потока мощности на расстоянии $d$ от передатчика равна:
\begin{equation}
    P = \frac{P_t}{4\pi d^2},
\end{equation}
где $P_t$ --- излучаемая мощность. Если используется направленная антенна с коэффициентом усиления $G_t$, выражение принимает вид $P = P_t G_t / (4\pi d^2)$. 

С учетом эффективной площади приемной антенны $S_r = G_r \lambda^2 / 4\pi$, мощность принимаемого сигнала определяется классической формулой Фрииса:
\begin{equation}
    P_r = \frac{P_t G_t G_r \lambda^2}{(4\pi)^2 d^2}.
    \label{eq:friis}
\end{equation}
Из (\ref{eq:friis}) следует, что в свободном пространстве мощность убывает обратно пропорционально квадрату расстояния ($1/d^2$). Формула Фрииса наглядно демонстрирует, что для увеличения дальности связи можно либо повышать мощность передатчика, либо использовать направленные антенны. Однако на практике мощность ограничена требованиями биологической защиты, а расширение полосы частот --- стоимостью лицензий. Поэтому на первый план выходят методы обработки сигналов, учитывающие реальную структуру канала.

В реальных условиях антенны подняты над землей, и в точку приема приходят как минимум два луча: прямой и отраженный от земной поверхности. В результате интерференции мощность принимаемого сигнала описывается двулучевой моделью:
\begin{equation}
    P_r = 4P_t G_t G_r \left(\frac{\lambda}{4\pi d}\right)^2 \sin^2\left(\frac{2\pi h_t h_r}{\lambda d}\right),
\end{equation}
где $h_t$ и $h_r$ --- высоты передающей и приемной антенн. Множитель $g = \sin^2(2\pi h_t h_r / \lambda d)$ называется множителем ослабления земли.

На расстояниях, превышающих граничное значение $d_{cp} = 4 h_t h_r / \lambda$, аргумент синуса становится малым, и мощность убывает обратно пропорционально четвертой степени расстояния:
\begin{equation}
    P_r \approx P_t G_t G_r \left(\frac{h_t h_r}{d^2}\right)^2.
\end{equation}
Эта закономерность известна как квадратичная формула Введенского. Важно отметить, что в данной зоне дальность связи перестает зависеть от несущей частоты. Зависимость множителя ослабления земли от расстояния показана на Рис.~\ref{fig:two_ray}.

\begin{figure}[htbp]
    \centering
    \begin{tikzpicture}
        \begin{axis}[
            width=12cm, height=5.5cm,
            xlabel={Расстояние $d$, км},
            ylabel={Множитель $g(d)$},
            xmin=0.1, xmax=10,
            ymin=0, ymax=1.1,
            xtick={0.1, 1, 10},
            xmode=log, 
            grid=both,
            draw=black,
            thick
        ]
            \addplot[domain=0.1:10, samples=500] 
                {sin(deg(2*pi*20*2/(0.33*x)))^2};
        \end{axis}
    \end{tikzpicture}
    \caption{Множитель ослабления земли}
    \label{fig:two_ray}
\end{figure}

\subsection{Крупномасштабные замирания}

Помимо регулярного ослабления, сигнал подвержен замираниям. Крупномасштабные (медленные) замирания возникают из-за затенения сигнала крупными объектами: зданиями, холмами, лесными массивами. Характерный радиус корреляции таких замираний составляет 20--50 метров.

Экспериментально установлено, что амплитуда сигнала при крупномасштабных замираниях подчиняется логарифмически нормальному закону распределения:
\begin{equation}
    p(s) = \frac{1}{s\sqrt{2\pi}\sigma} \exp\left(-\frac{(\ln s - m)^2}{2\sigma^2}\right),
\end{equation}
где $m$ --- среднее значение логарифма амплитуды, $\sigma$ --- стандартное отклонение (обычно 6--10 дБ).

На практике для расчета уровня сигнала используется интерполяционная модель:
\begin{equation}
    L(d) = L(d_0) + 10\alpha \lg\left(\frac{d}{d_0}\right) + x_\sigma,
\end{equation}
где $L(d_0)$ --- ослабление на опорном расстоянии $d_0$, $\alpha$ --- показатель ослабления (около 4 для городских условий), $x_\sigma$ --- случайная гауссова величина с нулевым средним. Длительность крупномасштабных замираний обычно составляет несколько секунд, что сопоставимо с длительностью передачи пакетов данных. Поэтому в современных системах для борьбы с такими замираниями применяют помехоустойчивое кодирование и перемежение данных.

\subsection{Мелкомасштабные замирания и импульсная характеристика}

В отличие от проводных линий связи, где сигнал распространяется строго по кабелю, в беспроводном канале среда распространения является открытой. Сигнал отражается от зданий, автомобилей и других объектов, что приводит к появлению множества копий исходного сигнала, приходящих с разными задержками. Для математического описания этого явления используется понятие импульсной характеристики.

Мелкомасштабные (быстрые) замирания обусловлены многолучевым распространением. Сигнал приходит на приемник в виде суммы множества лучей с различными задержками $\tau_n$ и комплексными коэффициентами передачи $\alpha_n$. Отклик канала на дельта-импульс называется импульсной характеристикой (ИХ):
\begin{equation}
    h(\tau) = \sum_n \alpha_n \exp(-j 2\pi f_c \tau_n)\, \delta(\tau - \tau_n).
\end{equation}
Передаточная функция канала $H(f)$, являющаяся преобразованием Фурье от ИХ, показывает, как канал ослабляет различные частотные компоненты:
\begin{equation}
    H(f) = \sum_n \alpha_n \exp(-j 2\pi f \tau_n).
\end{equation}

Если полоса сигнала $W$ значительно меньше величины $1/\tau_2$ (где $\tau_2$ --- задержка второго луча), то все частоты испытывают одинаковые замирания. Такой канал называется \textit{плоским} (частотно-неселективным). В противном случае канал является \textit{частотно-селективным}, и для борьбы с искажениями требуются эквалайзеры.

\subsection{Временная дисперсия в канале}

Разброс задержек приходящих лучей характеризует временную дисперсию канала. Средняя задержка определяется как:
\begin{equation}
    \langle \tau \rangle = \frac{\sum_n \tau_n P(\tau_n)}{\sum_n P(\tau_n)},
\end{equation}
где $P(\tau_n) = \langle |\alpha_n|^2 \rangle$ --- средняя мощность $n$-го луча. Мерой разброса задержек является среднеквадратическое отклонение (временная дисперсия):
\begin{equation}
    \sigma_\tau = \sqrt{\langle \tau^2 \rangle - \langle \tau \rangle^2}.
\end{equation}

Временная дисперсия напрямую связана с полосой когерентности канала $\Delta f_H \approx 1/\sigma_\tau$. Если ширина спектра передаваемого сигнала превышает полосу когерентности, канал становится частотно-селективным. Для современных широкополосных систем это является типичной ситуацией, что требует применения специальных методов оценивания. Пример дискретного спектра мощности задержек сигнала, который часто используется при моделировании, показан на Рис.~\ref{fig:delay_profile}. Из рисунка видно, что имеется шесть лучей, задержки которых относительно первого луча составляют от 0.31 до 2.51 мкс.

\begin{figure}[htbp]
    \centering
    \begin{tikzpicture}
        \begin{axis}[
            width=12cm, height=6cm,
            xlabel={Задержка $\tau$, мкс},
            ylabel={Мощность луча, дБ},
            xmin=-0.2, xmax=3,
            ymin=-25, ymax=5,
            xtick={0, 0.5, 1, 1.5, 2, 2.5, 3},
            ytick={0, -10, -20},
            grid=both,
            ycomb,
            draw=black, 
            mark=*,          
            mark options={fill=black, draw=black, scale=1.2},
            thick
        ]
            \addplot[draw=none] coordinates {(-0.2, 0) (3, 0)};
            \addplot coordinates {
                (0, 0) (0.31, -1) (0.71, -9) (1.09, -10) (1.73, -15) (2.51, -20)
            };
        \end{axis}
    \end{tikzpicture}
    \caption{Профиль мощности задержек многолучевого канала}
    \label{fig:delay_profile}
\end{figure}

\subsection{Статистические свойства замираний во времени}

Мелкомасштабные замирания представляют собой случайный процесс. В случае отсутствия прямого луча результирующий сигнал представляет собой сумму большого числа независимых отраженных сигналов. Согласно центральной предельной теореме, квадратурные компоненты $I(t)$ и $Q(t)$ такого сигнала подчиняются нормальному закону распределения с нулевым средним. 

Огибающая сигнала $A = \sqrt{I^2 + Q^2}$ в этом случае распределена по закону Релея:
\begin{equation}
    p(A) = \frac{A}{\sigma^2} \exp\left(-\frac{A^2}{2\sigma^2}\right) \quad (A > 0).
\end{equation}
Такой канал называется \textit{релеевским}. Фаза сигнала при этом распределена равномерно, а амплитуда и фаза статистически независимы.

Если же в точке приема присутствует детерминированная (регулярная) составляющая (например, прямой луч), распределение огибающей описывается распределением Райса:
\begin{equation}
    p(A) = \frac{A}{\sigma^2} \exp\left(-\frac{A^2 + A_0^2}{2\sigma^2}\right) I_0\left(\frac{A A_0}{\sigma^2}\right),
\end{equation}
где $A_0$ --- амплитуда регулярной компоненты, $I_0(\cdot)$ --- функция Бесселя первого рода нулевого порядка. Райсовский канал характеризуется K-фактором.

\subsection{Спектральные и корреляционные свойства}

Для описания того, как быстро меняются замирания во времени, используется функция корреляции. Из-за движения абонента возникает доплеровский сдвиг частоты $f_d = (v/\lambda)\cos\theta$. При равномерном распределении переотражателей вокруг приемника спектральная плотность мощности флуктуаций имеет вид доплеровского спектра Джейкса:
\begin{equation}
    \Phi_{zz}(f_d) = \frac{1}{\pi \sqrt{f_{\max}^2 - f_d^2}},
\end{equation}
где $f_{\max} = v/\lambda$ --- максимальное доплеровское смещение. Функция корреляции при этом определяется функцией Бесселя нулевого порядка:
\begin{equation}
    \phi_{zz}(\tau) = J_0(2\pi f_{\max} \tau).
\end{equation}
Величина $\tau_{coh} \approx 1/f_{\max}$ называется временем когерентности канала. Если длительность символа меньше времени когерентности, замирания считаются медленными (квазистатическими).

\subsection{Угловая и пространственная дисперсия}

В многолучевом канале сигналы приходят не только с разными задержками, но и с разных направлений. Это явление называется угловой дисперсией. Для базовой станции источник сигнала видится протяженным в пределах сектора $\theta_{\text{eff}}$. Угловая дисперсия влияет на пространственную корреляцию сигналов в разнесенных антеннах:
\begin{equation}
    \rho_{12} = J_0\left(\frac{2\pi d}{\lambda}\right).
\end{equation}
Из этого следует важный практический вывод: для обеспечения некоррелированного приема (для MIMO-систем) антенны нужно разносить на расстояние порядка половины длины волны ($d \approx 0.5\lambda$).

\subsection*{Выводы по разделу}

Таким образом, беспроводной канал связи представляет собой сложную случайную систему. Он характеризуется регулярным ослаблением сигнала с расстоянием, медленными замираниями из-за затенения и быстрыми замираниями из-за многолучевости. Знание импульсной характеристики канала и ее статистических свойств является основой для построения эффективных алгоритмов оценивания информационных символов и эквализации.